% Level 2, Week 12 answers -- Showcase and Celebration
\answersweek{12}{Showcase and Celebration}{Answers and scoring notes for workbook chapter 12 and homework 12.}

\begin{teachernote}[Exit assessment]
The answers and rubric for the Exit assessment in the Assessments booklet are in the assessments
section at the end of this answer key. Score it after the session, not during the celebration.
\end{teachernote}

\section{Workbook exercises}
\begin{answer}[Exercise 12.1: My one-minute pitch]
A strong pitch names the project, explains one real part of the code (a function with what goes in
and what comes out, or the loop, list, dictionary or file), gives one reason for a design choice,
and says what is next. Example: \enquote{Our \texttt{load\_high\_score} function opens
\texttt{highscore.txt} and returns the number in it, or 0 if the file is not there yet.} Use the pitch
as evidence for project goal L2.G8 (explaining design).
\end{answer}

\begin{answer}[Exercise 12.2: Gallery walk feedback]
Good stars point at something specific (\enquote{the rocks are squares, so I can see them without
color}); good wishes are kind and doable (\enquote{a slower first level would help me}). If a card
says only \enquote{cool}, ask the writer \enquote{which part exactly?} and let them add a word.
\end{answer}

\begin{answer}[Exercise 12.3: My Level 2 skills check]
There are no right answers. Compare it with the exit assessment: if a student says \enquote{not yet}
but showed the skill, tell them! Lines marked \enquote{not yet} are good starting points for the
holiday menu or a buffer session before Level 3 (the Level 3 placement quiz covers Level 2).
\end{answer}

\begin{answer}[Exercise 12.4: Looking ahead to Level 3]
Any three features. Nudge toward ideas that need a new data structure or a new test, for example
\enquote{a top-5 high-score table} (list in a file) or \enquote{difficulty levels} (a dictionary of
settings per level).
\end{answer}

\section{Homework}
\begin{answer}[Core: My Level 2 reflection]
Look for specific answers (a named bug and how it was fixed, a named piece of feedback). Vague
answers (\enquote{everything}) are fine from a tired student; you might ask one follow-up question
in person instead of correcting the paper.
\end{answer}

\begin{answer}[Challenge: Version 2 plan]
Example: \enquote{Feature: lives. Function \texttt{lose\_life(lives)} → \texttt{lives - 1} but never below
0. Test: \texttt{assert lose\_life(0) == 0}.} Each feature should have a function, its input and
output, and one test idea.
\end{answer}

\begin{answer}[Creative: holiday menu]
Optional. Celebrate anything they bring back in Level 3 week 1, and check that anything shared used
nicknames only.
\end{answer}
